\begin{exercises}
  \recommended \item
    对来自\nearbyexample{ex:DiagTwoByTwo}的矩阵，将\nearbyexample{ex:DiagUpperTrian}重做一遍。
    \begin{answer}
      由于基向量是任意选取的，因此可能有许多不同的答案。
      不过，下面给出一种做法；
      要对角化
      \begin{equation*}
         T=\begin{mat}[r]
            4  &-2  \\
            1  &1
         \end{mat}
      \end{equation*}
      把它视为某个变换在标准基下的表示 $T=\rep{t}{\stdbasis_2,\stdbasis_2}$，并寻找
      \( B=\sequence{\vec{\beta}_1,\vec{\beta}_2} \) 使得
      \begin{equation*}
        \rep{t}{B,B}
        =
        \begin{mat}
          \lambda_1  &0          \\
          0          &\lambda_2
        \end{mat}
      \end{equation*}
      也就是说，使得
      $t(\vec{\beta}_1)=\lambda_1$ 且 $t(\vec{\beta}_2)=\lambda_2$。
      \begin{equation*}
        \begin{mat}[r]
           4  &-2  \\
           1  &1
        \end{mat}
        \vec{\beta}_1=\lambda_1\cdot\vec{\beta}_1
        \qquad
        \begin{mat}[r]
           4  &-2  \\
           1  &1
        \end{mat}
        \vec{\beta}_2=\lambda_2\cdot\vec{\beta}_2
      \end{equation*}
      我们要寻找标量 \( x \)，使得方程
      \begin{equation*}
       \begin{mat}[r]
          4  &-2  \\
          1  &1
       \end{mat}
       \colvec{b_1 \\ b_2}=x\cdot\colvec{b_1 \\ b_2}
      \end{equation*}
      有不全为零的解 $b_1$ 与 $b_2$。
      将其改写为线性方程组
      \begin{equation*}
        \begin{linsys}{2}
           (4-x)\cdot b_1  &+  &-2\cdot b_2       &=  &0  \\
           1\cdot b_1      &+   &(1-x)\cdot b_2   &=  &0
        \end{linsys}
      \end{equation*}
      若 $x=4$，则第一个方程给出 $b_2=0$，进而第二个方程给出 $b_1=0$。
      我们已排除两个 $b$ 都为零的情形，
      故可假设 $x\neq 4$。
      \begin{equation*}
        \grstep{(-1/(4-x))\rho_1+\rho_2}
        \begin{linsys}{2}
           (4-x)\cdot b_1  &+   &-2\cdot b_2                   &=  &0  \\
                           &    &((x^2-5x+6)/(4-x))\cdot b_2   &=  &0
        \end{linsys}
      \end{equation*}
      考虑最下面的方程。
      若 \( b_2=0 \)，则第一个方程给出 $b_1=0$ 或 $x=4$。
      $b_1=b_2=0$ 的情形是不允许的。
      最下面方程的另一种可能是分式 $x^2-5x+6=(x-2)(x-3)$ 的分子为零。
      $x=2$ 的情形给出第一个方程 $2b_1-2b_2=0$，因此与 $x=2$ 相关联的是
      第一、第二分量相等的向量：
      \begin{equation*}
         \vec{\beta}_1=\colvec[r]{1 \\ 1}
         \qquad\text{（于是 \(
           \begin{mat}[r]
              4  &-2  \\
              1  &1
           \end{mat}
           \colvec[r]{1 \\ 1}=2\cdot\colvec[r]{1 \\ 1} \)，且 $\lambda_1=2$）。}
      \end{equation*}
      若 \( x=3 \)，则第一个方程为
      $b_1-2b_2=0$，于是相关的向量
      是那些第一分量为第二分量两倍的向量：
      \begin{equation*}
         \vec{\beta}_2=\colvec[r]{2 \\ 1}
         \qquad\text{（于是 \(
           \begin{mat}[r]
              4  &-2  \\
              1  &1
           \end{mat}
           \colvec[r]{2 \\ 1}=3\cdot\colvec[r]{2 \\ 1} \)，从而 $\lambda_2=3$）。}
      \end{equation*}
      下图
        \begin{equation*}
          \begin{CD}
            \C^2_{\wrt{\stdbasis_2}}
               @>t>T>
               \C^2_{\wrt{\stdbasis_2}}       \\
            @V{\scriptstyle\identity} VV
               @V{\scriptstyle\identity} VV \\
            \C^2_{\wrt{B}}
               @>t>D>
               \C^2_{\wrt{B}}
          \end{CD}
        \end{equation*}
      展示了如何得到该对角化。
      \begin{equation*}
         \begin{mat}[r]
           2  &0  \\
           0  &3
         \end{mat}
         =
         \begin{mat}[r]
           1  &2  \\
           1  &1
         \end{mat}^{-1}
         \begin{mat}[r]
           4  &-2  \\
           1  &1
         \end{mat}
         \begin{mat}[r]
           1  &2  \\
           1  &1
         \end{mat}
      \end{equation*}
      \textit{评注}。
      在下述改名之下，这个等式与 $T=PSP^{-1}$ 的定义相符。
      \begin{equation*}
        T=
         \begin{mat}[r]
           2  &0  \\
           0  &3
         \end{mat}
         \quad
         P=
         \begin{mat}[r]
           1  &2  \\
           1  &1
         \end{mat}^{-1}
         \quad
         P^{-1}=
         \begin{mat}[r]
           1  &2  \\
           1  &1
         \end{mat}
         \quad
         S=
         \begin{mat}[r]
           4  &-2  \\
           1  &1
         \end{mat}
      \end{equation*}
    \end{answer}
  \recommended \item \label{exer:DiagTheseA}
    按照\nearbyexample{ex:DiagUpperTrian}的步骤对角化下面的矩阵。
    \begin{equation*}
      \begin{mat}[r]
      1  &1  \\
      0  &0
      \end{mat}
    \end{equation*}
    \begin{exparts}
      \partsitem 写出\nearbylemma{lm:DiagIffBasisOfEigens}中描述的矩阵—向量方程，
         并将其改写为线性方程组。
      \partsitem 通过考虑该方程组的解，找出两个向量来构成一个基。
         （分别考虑 $x=0$ 与
         $x\neq 0$ 两种情形下的方程组。
         另外，回忆零向量不能是基的成员。）
      \partsitem 在相似图中使用该基，把对角矩阵表示为三个矩阵的乘积。
    \end{exparts}
    \begin{answer}
     \begin{exparts}
        \partsitem 我们要一个基，使其中的向量满足
          \begin{equation*}
            \begin{mat}
              1 &1 \\
              0 &0
            \end{mat}
            \colvec{b_1 \\ b_2}
            =x\cdot \colvec{b_1 \\ b_2}
          \end{equation*}
          于是考虑由此得到的线性方程组。
          \begin{equation*}
            \begin{linsys}{2}
              b_1 &+ &b_2 &= &x\cdot b_1 \\
                  &  &0   &= &x\cdot b_2 \\
            \end{linsys}
          \end{equation*}
        \partsitem
          若 $x=0$，则 $b_1+b_2=0$，即 $b_1=-b_2$。
          也就是说，与这种情形相关联的基元素，
          其分量符号相反。
          （分量必须非零，因为零向量
          不能是基的成员。）

          若 $x\neq 0$，则 $b_2=0$，把它代入上面的
          方程得到 $b_1=x\cdot b_1$。
          所以在这种情形，要么 $b_1=0$，要么 $x=1$。
          但由于此情形中 $b_2=0$，$b_1=0$ 是不可能的，因为
          零向量不在任何基中。
          于是~$x=1$，方程组变为
          \begin{equation*}
            \begin{linsys}{2}
              b_1 &+ &b_2 &= &1\cdot b_1 \\
                  &  &0   &= &1\cdot b_2 \\
            \end{linsys}
          \end{equation*}
          所以这种情形相关联的基元素满足 $b_2=0$，
          而~$b_1$ 可取任意非零值。

          因此，合适的基是
          \begin{equation*}
            B=\sequence{\colvec{1 \\ -1},
                        \colvec{1 \\ 0}}
          \end{equation*}
        \partsitem
          有了这个基，
          为得到对角化，画出相似图。
          \begin{equation*}
            \begin{CD}
            \C^2_{\wrt{\stdbasis_2}}
               @>t>T>
               \C^2_{\wrt{\stdbasis_2}}       \\
            @V{\scriptstyle\identity} VV
               @V{\scriptstyle\identity} VV \\
            \C^2_{\wrt{B}}
               @>t>D>
               \C^2_{\wrt{B}}
           \end{CD}
          \end{equation*}
          在此图中，
          \begin{equation*}
             T=
             \begin{mat}
               1  &1  \\
               0  &0
             \end{mat}
          \end{equation*}
          我们将利用基~$B$，
          沿着左下到左上、再到右上、
          再到右下的路径来计算矩阵~$D$。

          我们先要从左下走到左上。
          注意到由于我们把图上方的基
          选为标准基 $\stdbasis_2$，
          表示该映射的矩阵很容易写出。
          \begin{equation*}
            \rep{\identity}{B,\stdbasis_2}=
            \begin{mat}
              1 &1 \\
             -1 &0
            \end{mat}
          \end{equation*}
          整条路径的计算确实给出了
          对角的结果
          （当然，左下到左上这一段出现在
          三个矩阵乘积的右边）。
          \begin{align*}
                 D
                 &=
                 \begin{mat}
                    1  &1  \\
                    -1  &0
                 \end{mat}^{-1}
                 \begin{mat}
                    1  &1  \\
                    0  &0
                 \end{mat}
                 \begin{mat}
                    1  &1  \\
                    -1  &0
                 \end{mat}                \\
                 &=
                 \begin{mat}
                    0  &-1  \\
                    1  &1
                 \end{mat}
                 \begin{mat}
                    1  &1  \\
                    0  &0
                 \end{mat}
                 \begin{mat}
                    1  &1  \\
                    -1  &0
                 \end{mat}
                =\begin{mat}
                    0  &0  \\
                    0  &1
                 \end{mat}
          \end{align*}
     \end{exparts}
    \end{answer}
  \recommended \item \label{exer:DiagTheseB}
    按照\nearbyexample{ex:DiagUpperTrian}对角化下面的矩阵。
    \begin{equation*}
      \begin{mat}[r]
      0  &1  \\
      1  &0
      \end{mat}
    \end{equation*}
    \begin{exparts}
      \partsitem 写出矩阵—向量方程，
         并将其改写为线性方程组。
      \partsitem 通过考虑该方程组的解，找出两个向量来构成一个基。
         （分别考虑 $x=0$ 与
         $x\neq 0$ 两种情形。
         另外，回忆零向量不能是基的成员。）
      \partsitem 利用这个基和相似图，
        把对角化表示为三个矩阵的乘积。
    \end{exparts}
    \begin{answer}
     \begin{exparts}
        \partsitem
          基中的成员满足
          \begin{equation*}
            \begin{mat}
              0 &1 \\
              1 &0
            \end{mat}
            \colvec{b_1 \\ b_2}
            =x\cdot \colvec{b_1 \\ b_2}
          \end{equation*}
          因此由此得到线性方程组
          \begin{equation*}
            \begin{linsys}{2}
                   &  &b_2 &= &x\cdot b_1 \\
               b_1 &  &    &= &x\cdot b_2 \\
            \end{linsys}
            \qquad\Longrightarrow\qquad
            \begin{linsys}{2}
               -xb_1 &+ &b_2  &= &0 \\
                 b_1 &- &xb_2 &= &0 \\
            \end{linsys}
          \end{equation*}
        \partsitem
          若 $x=0$，则 $b_1=b_2=0$，但零向量不是任何
          基的成员，所以这种情形不可能发生。

          另一种情形是 $x\neq 0$。
          \begin{equation*}
            \begin{linsys}{2}
               -xb_1 &+ &b_2  &= &0 \\
                 b_1 &- &xb_2 &= &0 \\
            \end{linsys}
            \grstep{(1/x)\rho_1+\rho_2}
            \begin{linsys}{2}
               -xb_1 &+ &b_2           &= &0 \\
                     &  &(-x+(1/x))b_2 &= &0 \\
            \end{linsys}
          \end{equation*}
          关注最下面的方程。
          要么 $b_2=0$，要么 $-x+(1/x)=0$。
          若 $b_2=0$，则代入上面的方程给出
          $b_1=0$，因为 $x\neq 0$ 是这种情形的假设。
          但零向量不是任何基的成员，所以这不可能
          发生。

          于是只剩这种可能。
          \begin{equation*}
            \frac{1}{x}-x=0
            \qquad\Longrightarrow\qquad
            0=\frac{1-x^2}{x}=\frac{(1-x)(1+x)}{x}
          \end{equation*}
          若 $x=1$，则方程组为
          \begin{equation*}
            \begin{linsys}{2}
               -1\cdot b_1 &+ &b_2  &= &0 \\
                           &  &0    &= &0 \\
            \end{linsys}
          \end{equation*}
          于是基向量满足 $b_2=b_1$。
          若 $x=-1$，则方程组为
          \begin{equation*}
            \begin{linsys}{2}
               b_1 &+ &b_2  &= &0 \\
                   &  &0    &= &0 \\
            \end{linsys}
          \end{equation*}
          且基向量满足 $b_2=-b_1$。
          因此，下面是一个合适的基的例子。
          \begin{equation*}
            B=\sequence{\colvec{1 \\ 1},
                        \colvec{1 \\ -1}}
          \end{equation*}
        \partsitem
          为得到对角化，画出相似图
          \begin{equation*}
            \begin{CD}
            \C^2_{\wrt{\stdbasis_2}}
               @>t>T>
               \C^2_{\wrt{\stdbasis_2}}       \\
            @V{\scriptstyle\identity} VV
               @V{\scriptstyle\identity} VV \\
            \C^2_{\wrt{B}}
               @>t>D>
               \C^2_{\wrt{B}}
           \end{CD}
          \end{equation*}
          其中 $T$ 是题目中给出的矩阵。
          \begin{equation*}
             T=
             \begin{mat}
               0  &1  \\
               1  &0
             \end{mat}
          \end{equation*}
          利用基~$B$ 计算矩阵~$D$，
          从左下走到左上，再到右上，
          然后到右下。

          注意表示左下到左上
          映射的矩阵很容易写出。
          \begin{equation*}
            \rep{\identity}{B,\stdbasis_2}=
            \begin{mat}
              1 &1 \\
              1 &-1
            \end{mat}
          \end{equation*}
          有了它，整条路径的计算确实给出
          对角的结果。
          \begin{align*}
                 D
                 &=
                 \begin{mat}
                    1  &1  \\
                    1  &-1
                 \end{mat}^{-1}
                 \begin{mat}
                    0  &1  \\
                    1  &0
                 \end{mat}
                 \begin{mat}
                    1  &1  \\
                    1  &-1
                 \end{mat}                \\
                 &=
                 -\frac{1}{2}
                 \begin{mat}
                    -1  &-1  \\
                    -1  &1
                 \end{mat}
                 \begin{mat}
                    0  &1  \\
                    1  &0
                 \end{mat}
                 \begin{mat}
                    1  &1  \\
                    1  &-1
                 \end{mat}
                =\begin{mat}
                    1  &0  \\
                    0  &-1
                 \end{mat}
          \end{align*}
      \end{exparts}
    \end{answer}
  \item
    对角化这些上三角矩阵。
    \begin{exparts*}
      \partsitem
        $\begin{mat}[r]
          -2  &1  \\
           0  &2
        \end{mat}$
      \partsitem
        $\begin{mat}[r]
          5  &4  \\
          0  &1
        \end{mat}$
    \end{exparts*}
    \begin{answer}
      \begin{exparts}
        \partsitem
          写出
          \begin{equation*}
            \begin{mat}[r]
              -2  &1  \\
               0  &2
            \end{mat}
            \colvec{b_1 \\ b_2}
            =x\cdot\colvec{b_1 \\ b_2}
            \qquad\Longrightarrow\qquad
            \begin{linsys}{2}
               (-2-x)\cdot b_1  &+  &b_2            &=  &0  \\
                                &   &(2-x)\cdot b_2 &= &0
            \end{linsys}
          \end{equation*}
          给出两种可能：$b_2=0$ 与 $x=2$。
          沿 $b_2=0$ 这种可能，第一个方程为
          $(-2-x)b_1=0$，有两种情形：$b_1=0$ 或
          $x=-2$。
          于是，在第一种可能之下，我们得到 $x=-2$ 以及
          相关联的向量，其第二分量为零，而
          第一分量自由。
          \begin{equation*}
            \begin{mat}[r]
              -2  &1  \\
               0  &2
            \end{mat}
            \colvec{b_1 \\ 0}
            =-2\cdot\colvec{b_1 \\ 0}
            \qquad
            \vec{\beta}_1=\colvec[r]{1 \\ 0}
          \end{equation*}
          沿另一种可能，得到第一个方程
          $-4b_1+b_2=0$，因此与这个
          解相关联的向量的第二分量是其第一分量
          的四倍。
          \begin{equation*}
            \begin{mat}[r]
              -2  &1  \\
               0  &2
            \end{mat}
            \colvec{b_1 \\ 4b_1}
            =2\cdot\colvec{b_1 \\ 4b_1}
            \qquad
            \vec{\beta}_2=\colvec{1 \\ 4}
          \end{equation*}
          对角化如下。
          \begin{equation*}
            \begin{mat}[r]
              1  &1  \\
              0  &4
            \end{mat}
            \begin{mat}[r]
              -2  &1  \\
               0  &2
            \end{mat}
            \begin{mat}[r]
              1  &1  \\
              0  &4
            \end{mat}^{-1}
            =
            \begin{mat}[r]
              -2  &0  \\
               0  &2
            \end{mat}
          \end{equation*}
      \partsitem 计算与前一小题类似。
          \begin{equation*}
            \begin{mat}[r]
               5  &4  \\
               0  &1
            \end{mat}
            \colvec{b_1 \\ b_2}
            =x\cdot\colvec{b_1 \\ b_2}
            \qquad\Longrightarrow\qquad
            \begin{linsys}{2}
               (5-x)\cdot b_1  &+  &4\cdot b_2     &=  &0  \\
                               &   &(1-x)\cdot b_2 &=  &0
            \end{linsys}
          \end{equation*}
          最下面的方程
          给出两种可能：$b_2=0$ 与 $x=1$。
          沿 $b_2=0$ 这种可能，并舍去
          $b_2$ 与 $b_1$ 都为零的情形，得到
          $x=5$，与之相关联的向量第二分量
          为零，而第一分量自由。
          \begin{equation*}
            \vec{\beta}_1=\colvec[r]{1 \\ 0}
          \end{equation*}
          $x=1$ 这种可能给出第一个方程
          $4b_1+4b_2=0$，因此相关联的向量的
          第二分量是其第一分量的相反数。
          \begin{equation*}
            \vec{\beta}_1=\colvec[r]{1 \\ -1}
          \end{equation*}
          于是我们有如下对角化。
          \begin{equation*}
            \begin{mat}[r]
              1  &1  \\
              0  &-1
            \end{mat}
            \begin{mat}[r]
               5  &4  \\
               0  &1
            \end{mat}
            \begin{mat}[r]
              1  &1  \\
              0  &-1
            \end{mat}^{-1}
            =
            \begin{mat}[r]
               5  &0  \\
               0  &1
            \end{mat}
          \end{equation*}
      \end{exparts}
    \end{answer}
  \item 若试图用
    \nearbyexample{ex:DiagUpperTrian}的方法
    对角化\nearbyexample{ex:MatrixNotDiagonalizable}中的矩阵
    \begin{equation*}
       N=\begin{mat}[r]
            0  &0  \\
            1  &0
         \end{mat}
    \end{equation*}
    那么会在哪里出问题？
    \begin{exparts}
      \partsitem 画出带 $N$ 的相似图。
      \partsitem 写出\nearbylemma{lm:DiagIffBasisOfEigens}中描述的矩阵—向量方程，
         并将其改写为线性方程组。
      \partsitem 通过考虑该方程组的解，找出
         问题所在。
         （分别考虑 $x=0$ 与
         $x\neq 0$ 两种情形。）
    \end{exparts}
    \begin{answer}
      \begin{exparts}
        \item  在这个相似图
          \begin{equation*}
            \begin{CD}
            \C^2_{\wrt{\stdbasis_2}}
               @>n>N>
               \C^2_{\wrt{\stdbasis_2}}       \\
            @V{\scriptstyle\identity} VV
               @V{\scriptstyle\identity} VV \\
            \C^2_{\wrt{B}}
               @>n>D>
               \C^2_{\wrt{B}}
           \end{CD}
          \end{equation*}
          中，我们将尝试计算矩阵~$D$ 并找出合适的
          基~$B$，
          即沿左下到左上、再到右上、
          再到右下的路径行进。
        \item 该图表明，当
          $B=\sequence{\vec{\beta}_1,\vec{\beta}_2}$ 时，我们要
          $n(\vec{\beta}_1)=\lambda_1\vec{\beta}_1$ 且
          $n(\vec{\beta}_2)=\lambda_2\vec{\beta}_2$。
          为确定 $\lambda_1$ 和 $\lambda_2$，
          由于矩阵~$N$ 表示的是关于
          标准基的映射，我们要求这个方程组的解。
          \begin{equation*}
            \begin{mat}
              0 &0 \\
              1 &0
            \end{mat}
            \colvec{b_1 \\ b_2}
            =
            x\cdot\colvec{b_1 \\ b_2}
          \end{equation*}
          同解的线性方程组为
          \begin{equation*}
            \begin{linsys}{2}
              0b_1 &+  &0b_2 &= &xb_1 \\
              b_1  &+  &0b_2 &= &xb_2 \\
            \end{linsys}
            \qquad
            \Longrightarrow
            \qquad
            \begin{linsys}{2}
              -xb_1 &   &     &= &0 \\
              b_1   &-  &xb_2 &= &0 \\
            \end{linsys}
          \end{equation*}
          由第一行，要么 $x=0$，要么~$b_1=0$。
          若 $x=0$，则第二个方程变为 $b_1-0=0$，故 $b_1=0$，
          于是得到这个解集。
          \begin{equation*}
            \set{\colvec{0 \\ 1}b_2 \suchthat b_2\in\C}
          \end{equation*}
          我们可以取基的第一个成员为
          $\vec{\beta}_1=\vec{e}_2$。

          若 $x\neq 0$，则线性方程组的
          第一个方程给出 $b_1=0$。
          方程组的第一行变为 $0=0$，第二行变为
          $-xb_2=0$。
          由于这种情形取 $x\neq 0$，我们得到 $b_2=0$，
          解集是平凡的。
          \begin{equation*}
            \set{\colvec{0 \\ 0}}
          \end{equation*}
          所以出问题的地方在于：我们得不到
          \nearbylemma{lm:DiagIffBasisOfEigens}所要求的基。
      \end{exparts}
    \end{answer}
